\indent ಹೀಗೆ, ಸಂಧಿಕಾರ್ಯವು ನಡೆಯುವಾಗ ಒಂದು ಅಪ್ರ\par

\kbvsechead{ಬದಲು ಮತ್ತೊಂದು ಅಕ್ಷರ ಬಂದು ಸೇರುವುದಕ್ಕೆ “ಆದೇಶ ಸಂ:}

\noindent\hspace*{2.2em}ಎಂದು ಹೆಸರು.\par

\vspace{3pt plus 2pt}\noindent\hspace*{1.2em}\kbvlabel{ಉದಾ}\par\vspace{1pt}

\begin{kbvderiv}
\vspace{4pt plus 3pt minus 1pt}
\begin{center}\begin{adjustbox}{max width=\textwidth}
{\renewcommand{\arraystretch}{1.22}\setlength{\tabcolsep}{4pt}
\begin{tabular}{l c l}
ಊರು + ಕೋಲು & {\color{kbvgrey}$=$} & \kbvresult{ಊರುಗೋಲು} \\
ಕೆಳ + ತುಟಿ & {\color{kbvgrey}$=$} & \kbvresult{ಕೆಳದುಟಿ.} \\
\end{tabular}}
\end{adjustbox}\end{center}
\vspace{4pt plus 3pt minus 1pt}
\end{kbvderiv}

\indent ಈ ಉದಾಹರಣೆಗಳಲ್ಲಿ ಕ್ರಮವಾಗಿ ಕ-ಕಾರಕ್ಕೆ ಬದಲಾಗಿ ಗ-ಕಾರವೂ, ತ-ಕಾರಕ್ಕೆ ಬದಲಾಗಿ ದ-ಕಾರವೂ, ಪ-ಕಾರಕ್ಕೆ ಬದಲಾಗಿ ಬ-ಕಾರವೂ ಆದೇಶವಾಗಿ ಬಂದಿರುವುದನ್ನು ಗಮನಿಸಿ ; ಇವುಗಳೆಲ್ಲವೂ ವ್ಯಂಜನಗಳು.\par

\kbvsechead{ನೆನಪಿಡಿ :}

\vspace{2.5pt plus 3pt minus 0.5pt}\noindent\hspace*{1.2em}\kbvnum{(೧)}\hspace{0.55em}ಸ್ವರಾಕ್ಷರದ ಮುಂಜಿ ಸ್ವರಾಕ್ಷರವು ಬಂದು ಸಂಧಿಕಾರ್ಯ ವುಂಟಾದಲ್ಲಿ, ಅದು ಸ್ವರ ಸಂಧಿ ಎನಿಸುವುದು. ಲೋಪ ಮತ್ತು ಆಗಮ ಸಂಧಿಗಳು ಸ್ವರ ಸಂಧಿಗಳು.\par

\vspace{2.5pt plus 3pt minus 0.5pt}\noindent\hspace*{1.2em}\kbvnum{(೨)}\hspace{0.55em}ಸ್ವರಾಕ್ಷರದ ಮುಂಜಿ ವ್ಯಂಜನವು ಬಂದು ಸಂಧಿಕಾರ್ಯ ವುಂಟಾದಲ್ಲಿ, ಅದು ವ್ಯಂಜನ ಸಂಧಿ ಎನಿಸುವುದು. ಆದೇಶ ಸಂಧಿಯು ವ್ಯಂಜನ ಸಂಧಿಯಾಗಿರುವುದು.\par

\vspace{2.5pt plus 3pt minus 0.5pt}\noindent\hspace*{1.2em}\kbvnum{(೩)}\hspace{0.55em}ಸ್ವರದ ಮುಂಜಿ ಸ್ವರವು ಬಂದರೂ ಸಂಧಿಕಾರ್ಯವಾಗದಿರುವ ಸಂದರ್ಭಗಳುಂಟು, ಇದನ್ನು ಪ್ರಕೃತಿಭಾನ ಎನ್ನುವರು. (ಸಂಧಿಕಾರ್ಯ ವಾಗಬೇಕಾದ ಸಂದರ್ಭಗಳಲ್ಲಿ ಪದಗಳು ಇದ್ದ ಹಾಗೆಯೇ ಇರುವುದೇ ಪ್ರಕೃತಿಭಾನ.)\par

\vspace{3pt plus 2pt}\noindent\hspace*{1.2em}\kbvlabel{ಉದಾ}\par\vspace{1pt}

\begin{kbvderiv}
\vspace{4pt plus 3pt minus 1pt}
\noindent{\renewcommand{\arraystretch}{1.24}\setlength{\tabcolsep}{5pt}
\begin{tabularx}{\textwidth}{X X X X}
ಆ ಅಗಸ & ಎಲಾ ಅದೇನು & ಬಾಣ ಇಲ್ಲಿ & ಮುದಿ ಆನೆ \\
ಲಕ್ಷ್ಮಣಾ ! ಅತ್ತ ನೋಡು & ಅಬ್ಬಾ ! ಅದೇನು ಸಂಭ್ರಮ &  &  \\
\end{tabularx}}
\vspace{4pt plus 3pt minus 1pt}
\end{kbvderiv}
